Data assimilation (DA) estimates unknown states or parameters in dynamical systems by combining prior model predictions with noisy observations, typically through variational methods or Bayesian methods that estimate posterior distributions via finite-sample approximations \cite{evensen_data_2022}. The Bayesian approach is often preferred for its natural uncertainty quantification. Within Bayesian DA, ensemble Kalman filters are efficient but require Gaussianity assumptions and handle only weakly nonlinear models, while particle filters impose virtually no restrictions but are expensive due to the large number of samples required. This trade-off has motivated alternatives such as neural-network surrogates for expensive forward models \cite{mucke_deep_2024, cheng_generalised_2022, zhong_reduced-order_2023, chen_reduced-order_2023}; these deterministic approximations require specifying a priori model error distributions, which is challenging for uncertainty quantification in high-dimensional stochastic systems.

Generative models have emerged as a compelling alternative, learning to sample directly from distributions of interest; in DA they are trained on prior distributions and modified to incorporate observational constraints. Two families of flow-based generative models dominate: stochastic interpolants (SIs) \cite{albergo_stochastic_2023, albergo_building_2023, chen_probabilistic_2024}, which transport samples through a stochastic differential equation (SDE), and flow matching (FM) \cite{lipman_flow_2023}, which does so through a deterministic ordinary differential equation (ODE). Both are increasingly applied to physical forecasting, yet methods for turning them into posterior samplers have so far been developed model by model, often tied to whether the generator is stochastic or deterministic.

\paragraph{A unified view.}
For data assimilation, the distinction between SDE- and ODE-based generators is a choice of sampler rather than a difference in method. Stochastic interpolants, flow matching, and the SDE form of flow matching -- which coincides with denoising diffusion -- are all special cases of a single \emph{affine Gaussian probability path} \cite{albergo_stochastic_2023, lipman_flow_2023, ma_sit_2024, holderrieth_generator_2024}. This unification is itself established \cite{albergo_stochastic_2023, holderrieth_generator_2024}; our contribution is its use as a single interface for posterior sampling, obtained by conditioning the path on the observation. Conditioning tilts the target to the posterior while leaving the path structure intact, so the conditional score and velocity differ from their prior counterparts by the \emph{same} likelihood-score term -- scaled by one for the score and by the velocity--score duality coefficient for the velocity \cite{albergo_stochastic_2023, ma_sit_2024}. Because the conditioned path again admits a one-parameter family of equal-marginal dynamics with a freely chosen diffusion coefficient \cite{song_score-based_2021}, a single guidance term yields a whole family of posterior samplers. Three are of practical interest: the SI-SDE (native diffusion), the FM-SDE (lifting the probability-flow ODE), and a deterministic guided probability-flow ODE that steers flow matching without injecting noise \cite{yan_fig_2024, pokle_training-free_2024}. All three share one observation-interpolant likelihood, differing only in a diffusion coefficient (Figure~\ref{fig:visual_abstract}). Treating these models as one family lets us state \emph{which} efficiency/accuracy trade-offs to make and \emph{when} they matter -- and when the samplers coincide so the cheapest choice suffices.

\input{figures/visual_abstract}

\paragraph{Contributions:}
We introduce a posterior sampling method that applies to any flow-based generative model defined by an affine Gaussian probability path, without retraining. Our key contributions:
\begin{itemize}
    \item A model-agnostic formulation that conditions the affine Gaussian path on the observation; the conditioned path is again affine Gaussian, its score and velocity correcting the prior ones by a single likelihood-score term, and we prove the marginal-preserving family generates exact posterior samples under ideal conditions.
    \item Three posterior samplers from this single family -- the SI-SDE, FM-SDE, and a deterministic guided ODE (native, lifted, and zero diffusion) -- sharing one guidance weight $\gweight_\tau=\vscoef_\tau+\tfrac12\gdiff_\tau^2$, with pseudo-code for each.
    \item A practical likelihood-score estimator from observation interpolation, its conditional moments for both model classes, and a closed-form interpolant likelihood whose covariance is inflated ($\Pi$GDM-style) by the model's source covariance \cite{chung_diffusion_2023, song_pseudoinverse_2023}; this inflation, shared by all three samplers, is what drives accuracy on sparse observations.
    \item An explicit account of which approximations control efficiency versus accuracy and when each matters (Tables~\ref{tab:samplers}--\ref{tab:si_vs_fm}), a relation to Doob's $h$-transform, and demonstrations on a series of test cases.
\end{itemize}
